\section{预训练优化}
\label{sec:nlp-llm-optimization}

目标函数给出了希望降低的误差，尚未决定多少参数、多少训练词元能够在既定预算内配合，也未决定怎样把梯度变为参数变化。下面固定前面的训练分布与目标，依次讨论预算配置、参数尺度、更新方向与训练日程。这些属于学习算法及其条件；训练设备怎样通信和分配计算不在本节范围内。

\subsection{模型规模、词元规模与计算预算}
\label{sec:nlp-llm-compute-allocation}

在固定训练计算量下，扩大模型通常意味着能够处理的训练词元减少。若参数多而材料曝光不足，模型可能尚未充分学习；若模型太小，即使继续增加词元，也可能受表达能力限制。因此，比较规模时应同时给出参数量$N$、训练词元量$D$及计算预算$C$，而不是只用其中一个数字解释结果。

对结构与序列设置相近的稠密模型，可以用$C\approx\kappa ND$概括训练计算量，其中$\kappa$吸收一次前向与反向的系数。这是便于比较的近似，长上下文的注意力开销、输出映射和不同结构都可能使系数变化。Chinchilla在多种参数和词元规模上拟合训练结果，说明仅扩大参数而不相应增加训练词元，可能偏离所研究预算范围内较好的配置\scite{hoffmann2022chinchilla}。

一种便于理解的经验损失形式为
\begin{equation}
 L(N,D)=L_{\infty}+A N^{-p}+B D^{-q},
 \qquad A,B,p,q>0.
 \label{eq:nlp-llm-scaling-loss}
\end{equation}
$L_{\infty}$表示拟合中的不可约部分，后两项描述有限模型和有限数据带来的额外损失。代入$D=C/(\kappa N)$后，对$N$求驻点，得到
\begin{equation}
 pA N^{-p}=qB D^{-q},\qquad
 N_{\mathrm{opt}}\propto C^{q/(p+q)},\quad
 D_{\mathrm{opt}}\propto C^{p/(p+q)}.
 \label{eq:nlp-llm-scaling-allocation}
\end{equation}
第一式表示，在这个拟合模型和预算约束下，两种规模的边际改善需要平衡。若$p$与$q$接近，增加预算时两者大致同阶增长；这不是由概率论推导出的普遍定律，而依赖损失拟合、数据分布和可比较的训练过程。

训练词元量还不等于独立新材料量。把十亿词元重复十遍虽产生一百亿次曝光，但后几遍仍是相同材料，不能无条件套用不断加入新数据时的拟合关系。若同一模型将在训练后被调用很多次，总预算还包含使用计算，可能愿意多训练较小模型以降低每次使用成本。稀疏专家模型又需要分别记录总参数和激活参数；用哪一个量代入经验式，必须重新建立相应比较，不能直接继承稠密模型系数。

\subsection{参数化、初始化与更新尺度}
\label{sec:nlp-llm-parameter-scale}

选定规模后，相同的学习率数值未必在不同宽度上产生相同大小的函数变化。对一层的简化初始化分析，设输入$\bm h$为$d$维，一行权重$\bm w$的元素零均值、互相独立、方差为$s^2$，并独立于输入。给定$\bm h$时，线性输出满足$\operatorname{Var}_{\bm w}(\bm w^{\mathsf T}\bm h\mid\bm h)=s^2\lVert\bm h\rVert_2^2$。若各输入分量的二阶矩均为$\nu$，再对输入平均后，输出方差就是$ds^2\nu$；这是上述初始化假设下的计算，不表示训练后权重和激活仍保持独立。若宽度增大却保持$s^2$不变，初始激活尺度就会变化；让$s^2$随$d$缩小可以控制初始尺度，但训练更新还需另行检查。

一次参数变化$\Delta\theta$对输出的局部影响近似为$\Delta f\approx J_\theta\Delta\theta$，其中$J_\theta$是输出对参数的雅可比矩阵。即使$\Delta\theta$中每个元素大小相同，宽度变化也会改变参与相加的项数及$J_\theta$的尺度。因此，初始化合适并不自动保证更新合适；嵌入、隐藏矩阵、输出映射的形状不同，也不能把同一缩放规则不加区分地应用到所有参数。

最大更新参数化（maximal update parametrization, $\mu$P）协调不同参数组的初始化与学习率尺度，使宽度增加时隐藏特征的更新保持在相应的尺度上。$\mu$Transfer在这类参数化条件下，探索把较小模型上选择的超参数迁移到较大宽度，减少每个规模都完整调参的需要\scite{nlp-llm-mutransfer2022}。它处理的是跨规模的参数与更新关系，不是式\eqref{eq:nlp-llm-scaling-allocation}中模型与数据预算的最优分配。

迁移仍需固定哪些结构与目标保持一致，并在目标规模检查激活、梯度和损失。若改变层深、专家结构、词表、批次或训练分布，原有宽度关系未必充分描述新设置。可迁移的超参数也不意味着不需要验证；它提供的是有条件的调参组织方式，而非对任意模型免除优化检查的保证。

\subsection{梯度方向与矩阵优化}
\label{sec:nlp-llm-optimizer}

将全部参数按坐标排列后，每批样本给出梯度$\bm g^{(s)}$，优化器需要把它变成第$s$次更新。采用解耦权重衰减的Adam（Adam with decoupled weight decay, AdamW）分别维护一阶与二阶逐坐标统计，用历史梯度的尺度调整当前步长。省略偏差校正的计算细节后，其核心形式为
\begin{align}
 \bm m^{(s)}&=\beta_1\bm m^{(s-1)}+(1-\beta_1)\bm g^{(s)},\\
 \bm v^{(s)}&=\beta_2\bm v^{(s-1)}+(1-\beta_2)(\bm g^{(s)})^{\odot2},\\
 \theta^{(s+1)}&=(1-\eta_s\lambda)\theta^{(s)}
  -\eta_s\frac{\widehat{\bm m}^{(s)}}
  {\sqrt{\widehat{\bm v}^{(s)}}+\varepsilon}.
 \label{eq:nlp-llm-adamw}
\end{align}
其中$\beta_1,\beta_2\in[0,1)$为平滑系数，帽号表示从零初始化统计量时的偏差校正，即$\widehat{\bm m}^{(s)}=\bm m^{(s)}/(1-\beta_1^s)$、$\widehat{\bm v}^{(s)}=\bm v^{(s)}/(1-\beta_2^s)$。除法和平方根逐坐标进行，$\eta_s$是学习率。$\lambda\geq0$控制与梯度预处理分开的权重衰减，$\varepsilon>0$避免零分母。哪些参数使用衰减，应在参数组中说明\scite{loshchilov2019adamw}。

矩阵权重除了包含许多坐标，还定义输入到输出的线性方向。Muon对选定矩阵参数的动量方向做近似正交化，再调整更新尺度。为解释这一步，设矩阵动量的紧奇异值分解为$\bm M=\bm U\bm\Sigma\bm V^{\top}$，可以用$\bm Q=\bm U\bm V^{\top}$代替原方向。奇异向量保留，非零奇异值的大小差别被移除；实际实现通常以矩阵迭代近似这一变换，不需要逐次求完整分解\scite{jordan2024muon}。

这种变换与AdamW的逐坐标重标度不同：前者根据矩阵的输入输出方向处理更新，后者根据各坐标的历史平方梯度处理更新。正交化也不意味着新权重矩阵本身被约束为正交矩阵，因为被处理的是更新方向，更新仍加到原有权重上。对于非方阵，还需要说明所取尺度与形状的关系；同一谱方向若整体放大不同倍，参数变化量也不同。

《Muon is Scalable for LLM Training》研究了将Muon用于大规模语言模型时的权重衰减与更新尺度等调整，隐藏矩阵与嵌入、输出或其他非矩阵参数可使用不同更新规则\scite{nlp-llm-muonscale2025}。该工作并非最初提出Muon，也不能据其特定训练比较断言它对所有结构和目标均优于AdamW。合理比较需要固定数据、预算、调参范围和评价口径；优化器名称本身不能代替这些条件。

\subsection{批次、学习率与训练日程}
\label{sec:nlp-llm-training-schedule}

单次梯度依赖抽到哪些样本。若每个目标词元都等权，一批含$B$个计分词元时，梯度可写为它们负对数损失梯度的平均；若实际采用按序列平均，则批次定义随之改变。在近似独立、方差有限的理想情况下，增加批次能降低平均梯度的抽样方差，但长文档相邻位置和重复样本并不独立，计分词元数不能直接当作相同数量的独立观测。

学习率日程规定$\eta_s$怎样随学习进展变化。预热阶段从较小步长逐渐增加，可以限制训练初期统计尚未稳定时的参数变化；之后的衰减减小后期更新幅度。一种示例是先线性预热至$\eta_{\max}$，再按余弦曲线降到$\eta_{\min}$。选择这些曲线的目的在于组织探索与细调，不是证明任意网络会收敛到相同解；还应说明日程横轴采用更新次数还是已处理词元量。

假设两个配置都训练一百万次，但每次计分词元数相差四倍，它们并没有使用相同数据曝光或计算预算。若批次随训练增大，又仍按更新次数衰减，就可能改变每个学习率阶段处理的材料数量。记录时应把批次、累计词元、学习率、梯度尺度与留出损失联系起来，才能区分是学习进展、数据变化还是异常更新造成曲线变化。

DeepSeek-V3等技术报告同时给出结构、主辅目标、优化超参数与多阶段日程，适合作为完整组合的实例\scite{deepseek2024v3}。单份报告中的成功组合不能分离每个选择的独立贡献；把其中一个超参数移到不同结构，需要重新检查。训练过程产出的参数、数据曝光和优化记录，正是下一节进行阶段检查与继续学习的依据。
